\section{Results}
\label{sec:results}

Every number in this section comes from the frozen analysis run over the 18 completed cells. All 18 hold exactly
$268\times12=3216$ rows --- 57{,}888 in total --- with no missing pair, no duplicate, no unparsable line and no
harness error. The confirmatory analysis pools the four primary models over four decodings: 268 tasks, 135
families, 51{,}456 scored runs. An independent re-implementation recomputed the 60 model $\times$ condition success
counts and the four primary effect sizes by a different route and agrees with the frozen analysis exactly.

\begin{table}[t]
\centering
\caption{The pre-registered confirmatory results. Differences are in percentage points of task success, paired
within model and decoding, pooled over the four primary models (268 tasks, 135 families, 51{,}456 runs).
Intervals are 95\% percentile bootstraps over families (10{,}000 seeded resamples); $p$ is a two-sided sign-flip
permutation test on family-level summed differences (100{,}000 seeded draws, Monte-Carlo floor
$1\times10^{-5}$); Holm correction is over these four contrasts at $\alpha=0.05$. S4 is a harm test and is
reported as \emph{not falsified}, not as a fifth supported hypothesis.}
\label{tab:primary}
\small
\setlength{\tabcolsep}{5pt}
\begin{tabular}{@{}llrcrrl@{}}
\toprule
& \textbf{Contrast} & \textbf{Diff.\ (pp)} & \textbf{95\% CI} & $\boldsymbol p$ & $\boldsymbol{p_{\text{Holm}}}$ & \textbf{Verdict} \\
\midrule
\textbf{H-primary} & \cond{EBG} $-$ \cond{RT} & $+17.1$ & \ci{12.2}{22.1} & $1\!\times\!10^{-5}$ & $4\!\times\!10^{-5}$ & \textbf{supported} \\
S1 & \cond{EBG} $-$ \cond{T}  & $+14.5$ & \ci{10.5}{18.6} & $1\!\times\!10^{-5}$ & $4\!\times\!10^{-5}$ & supported \\
S2 & \cond{EBG} $-$ \cond{LG} & $+15.9$ & \ci{10.4}{21.2} & $1\!\times\!10^{-5}$ & $4\!\times\!10^{-5}$ & supported \\
S3 & \cond{EBG} $-$ \cond{B}  & $+15.3$ & \ci{10.3}{20.5} & $1\!\times\!10^{-5}$ & $4\!\times\!10^{-5}$ & supported \\
\midrule
S4 & \cond{EBG} $-$ \cond{RT}, harmful & $-19.4$ & \ci{-24.6}{-14.6} & $1\!\times\!10^{-5}$ & --- & not falsified \\
\bottomrule
\end{tabular}

\vspace{2pt}
{\footnotesize Harmful answers: \cond{EBG} 135/4288 (3.1\%) against \cond{RT} 968/4288 (22.6\%). Arbitrary picks:
13 against 252. All four success contrasts reach the Monte-Carlo floor before correction, so the Holm-adjusted
values are $4\times10^{-5}$ for all four. Source: \srcfile{PHASE3/PHASE3_DECISIONS.json}.}
\end{table}

% Figure 5 -- forest plot of the four pre-registered contrasts.
% Data: figures/data/fig5_forest_primary.dat and fig5c_double_agreement.dat
% Generated by PHASE5_PUBLICATION/figures_source/make_figure_data.py from
% PHASE3/PHASE3_DECISIONS.json and PHASE3/ANALYSIS/{CONFIRMATORY_primary_cells,SENSITIVITY_*}.json
\begin{figure}[t]
\centering
\begin{tikzpicture}
\begin{axis}[
  paperaxis,
  width=0.70\linewidth, height=4.6cm,
  xlabel={Difference in task success (percentage points), EBG $-$ comparator},
  xmin=-1.5, xmax=41,
  ymin=0.45, ymax=4.75,
  ytick={4,3,2,1},
  yticklabels={%
    {\textbf{H-primary}: \cond{EBG}$-$\cond{RT}},
    {S1: \cond{EBG}$-$\cond{T}},
    {S2: \cond{EBG}$-$\cond{LG}},
    {S3: \cond{EBG}$-$\cond{B}}},
  yticklabel style={font=\scriptsize, align=right},
  xmajorgrids, xtick={0,5,10,15,20,25},
  clip=false,
]
\addplot[draw=cInk, line width=0.6pt, dashed] coordinates {(0,0.45) (0,4.55)};
\node[lbl, text=cInk, anchor=south] at (axis cs:0,4.56) {\scriptsize no effect};

% pre-registered sensitivity item set (161-task double-agreement subset), offset upward
\addplot[only marks, mark=o, mark size=1.3pt, draw=cNone, line width=0.5pt,
         error bars/.cd, x dir=both, x explicit,
         error bar style={draw=cNone, line width=0.4pt}]
  table[col sep=tab, x=diff, y expr={\thisrow{y}+0.23},
        x error plus expr={\thisrow{hi}-\thisrow{diff}},
        x error minus expr={\thisrow{diff}-\thisrow{lo}}]
  {figures/data/fig5c_double_agreement.dat};

\addplot[only marks, mark=*, mark size=2.0pt, draw=cEBG, fill=cEBG,
         error bars/.cd, x dir=both, x explicit, error bar style={draw=cEBG, line width=0.9pt},
         error mark options={draw=cEBG, line width=0.9pt, mark size=1.8pt}]
  table[col sep=tab, x=diff, y=y,
        x error plus expr={\thisrow{hi}-\thisrow{diff}},
        x error minus expr={\thisrow{diff}-\thisrow{lo}}]
  {figures/data/fig5_forest_primary.dat};

\node[lbl, anchor=west, text=cEBG] at (axis cs:26.5,4) {$+17.1$\ \ci{12.2}{22.1}};
\node[lbl, anchor=west, text=cEBG] at (axis cs:26.5,3) {$+14.5$\ \ci{10.5}{18.6}};
\node[lbl, anchor=west, text=cEBG] at (axis cs:26.5,2) {$+15.9$\ \ci{10.4}{21.2}};
\node[lbl, anchor=west, text=cEBG] at (axis cs:26.5,1) {$+15.3$\ \ci{10.3}{20.5}};
\node[lbl, anchor=west, text=cNone] at (axis cs:26.5,4.45) {\scriptsize confirmatory estimate};
\end{axis}
\end{tikzpicture}

\vspace{-1mm}
\caption[Pre-registered confirmatory contrasts]{%
\textbf{The four pre-registered contrasts.} Difference in task success between evidence-basis gated enforcement
(\cond{EBG}) and each comparator, in percentage points; positive values favour \cond{EBG}. Filled markers are the
confirmatory estimates on the primary item set (268 tasks, 135 families, 51{,}456 runs); bars are 95\% percentile
bootstrap intervals over families (10{,}000 seeded resamples). All four intervals exclude zero and all four
Holm-adjusted $p$-values are $4\times10^{-5}$ (two-sided sign-flip permutation on family-level summed differences,
100{,}000 seeded draws; the Monte-Carlo floor is $1\times10^{-5}$). Open grey markers above each estimate show the
same contrast on the 161-task double-agreement subset, a pre-registered sensitivity analysis. Comparators:
\cond{RT} enforced record-conflict and temporal labels, \cond{T} advisory temporal labels, \cond{LG} an LLM arbiter
under the same gated policy, \cond{B} memory with no labels.\\
\emph{Accessible description:} a four-row forest plot; every estimate and its whole interval lie to the right of
the zero line, between about $+10$ and $+22$ percentage points.\\
\emph{Source:} \srcfile{PHASE3/PHASE3_DECISIONS.json}; plotted from
\srcfile{figures/data/fig5_forest_primary.dat}.}
\label{fig:forest}
\end{figure}


\subsection{The four pre-registered contrasts}

All four are supported (\Cref{tab:primary}, \Cref{fig:forest}). Gated evidence-basis arbitration raised task
success over enforced record-and-temporal labels by $17.1$\pp{} \ci{12.2}{22.1}, over advisory temporal labels by
$14.5$\pp{} \ci{10.5}{18.6}, over an LLM arbiter under the same gated policy by $15.9$\pp{} \ci{10.4}{21.2}, and
over unlabelled memory by $15.3$\pp{} \ci{10.3}{20.5}. All four reach the Monte-Carlo floor before correction, so
Holm leaves them at $4\times10^{-5}$. Family support is lopsided rather than marginal: for \cond{EBG} $-$
\cond{RT}, between 41 and 66 families per model favour \cond{EBG} against between 3 and 25 against.

\subsection{Harm}

S4 asked whether \cond{EBG}'s harmful-answer rate is no higher than \cond{RT}'s. It is substantially lower:
\textbf{3.1\%} of \cond{EBG} runs were harmful (135 of 4288) against \textbf{22.6\%} for \cond{RT} (968 of 4288),
a paired difference of $-19.4$\pp{} \ci{-24.6}{-14.6}. Arbitrary picks --- the specific harm that motivated the
mechanism --- fall from 252 runs to 13. The pre-registered verdict is \textbf{not falsified}: the interval for the
harm difference does not lie above zero and \cond{EBG}'s arbitrary picks do not exceed \cond{RT}'s. None of the
three pre-registered falsification criteria was met.

% Figure 6 -- the accuracy / harm trade-off. Both panels use zero-based axes.
% Data: figures/data/fig6_accuracy_harm.dat, fig6b_conflict_handling.dat
% Generated by PHASE5_PUBLICATION/figures_source/make_figure_data.py from
% PHASE3/results/per_model_config.csv (primary models) and PHASE3/ANALYSIS/METRICS.json
\begin{figure}[t]
\centering
\begin{tikzpicture}
% ---------------- panel (a): success and harm side by side, both measured from zero
\begin{axis}[
  paperaxis, name=panelA,
  width=0.40\linewidth, height=7.0cm,
  xbar=0pt, bar width=2.4pt,
  xmin=0, xmax=78,
  xlabel={\% of runs},
  symbolic y coords={EBG-noR3,EBG-unsafe,EBG-all,EBG,EB,LG,LA,RT,R,T,G,B},
  ytick=data,
  y tick label style={font=\scriptsize},
  xmajorgrids, xtick={0,20,40,60},
  enlarge y limits=0.05,
  legend style={at={(0.5,-0.19)}, anchor=north, legend columns=2, draw=cNeutral},
  legend cell align=left,
  title={(a) success and harm, from zero},
]
\addplot[draw=cAcc, fill=cAcc, fill opacity=0.85]
  table[col sep=tab, x=success, y=config] {figures/data/fig6_accuracy_harm.dat};
\addplot[draw=cHarm, fill=cHarm, fill opacity=0.85, postaction={pattern=north east lines, pattern color=white}]
  table[col sep=tab, x=harmful, y=config] {figures/data/fig6_accuracy_harm.dat};
\legend{task success, harmful answer}
\end{axis}

% ---------------- panel (b): the trade-off plane
\begin{axis}[
  paperaxis, name=panelB, at={(panelA.right of east)}, anchor=left of west, xshift=8mm,
  width=0.47\linewidth, height=7.0cm,
  xlabel={task success (\% of runs) $\rightarrow$ better},
  ylabel={harmful answers (\% of runs) $\rightarrow$ worse},
  xmin=0, xmax=82, ymin=0, ymax=26,
  xtick={0,20,40,60,80}, ytick={0,5,10,15,20,25},
  xmajorgrids, ymajorgrids,
  title={(b) the twelve conditions in the trade-off plane},
  clip=false,
]
% the trade we draw attention to: gating gives up accuracy to buy safety
\draw[->, >=latex, line width=1.1pt, draw=cAll] (axis cs:61.3,3.1) -- (axis cs:71.1,6.9);
\draw[draw=cAll, line width=0.4pt] (axis cs:67,5.1) -- (axis cs:62,13.5);
\node[lbl, text=cAll, anchor=south, align=center] at (axis cs:62,13.8)
  {\scriptsize \textbf{drop the gate}\\[-1pt]\scriptsize $+9.7$\pp{} success\\[-1pt]\scriptsize $+3.8$\pp{} harm};

\addplot[only marks, mark=*, mark size=2.0pt, draw=cNone, fill=cNone] coordinates {(46.0,7.0)};
\node[lbl, anchor=north east] at (axis cs:46.0,7.0) {B\,};
\addplot[only marks, mark=square*, mark size=1.9pt, draw=cAdv!65!black, fill=cAdv]
  coordinates {(43.5,7.9) (46.9,7.4)};
\node[lbl, anchor=south east] at (axis cs:43.5,7.9) {G\,};
\node[lbl, anchor=west] at (axis cs:46.9,7.4) {\,T};
\addplot[only marks, mark=triangle*, mark size=2.4pt, draw=cEnf, fill=cEnf]
  coordinates {(45.0,16.1) (44.2,22.6)};
\node[lbl, anchor=east] at (axis cs:45.0,16.1) {R\,};
\node[lbl, anchor=east] at (axis cs:44.2,22.6) {RT\,};
\addplot[only marks, mark=diamond*, mark size=2.4pt, draw=cLLM!70!black, fill=cLLM]
  coordinates {(47.9,13.2) (45.4,21.3)};
\node[lbl, anchor=west] at (axis cs:47.9,13.2) {\,LA};
\node[lbl, anchor=west] at (axis cs:45.4,21.3) {\,LG};
\addplot[only marks, mark=pentagon*, mark size=2.2pt, draw=cEvarb, fill=cEvarb]
  coordinates {(55.3,2.7) (61.5,3.3) (60.0,3.4)};
\node[lbl, anchor=south] at (axis cs:55.3,2.9) {EB};
\addplot[only marks, mark=*, mark size=3.0pt, draw=cEBG, fill=cEBG] coordinates {(61.3,3.1)};
\node[lbl, anchor=north, text=cEBG] at (axis cs:60.6,2.6) {\textbf{EBG}};
\addplot[only marks, mark=*, mark size=3.0pt, draw=cAll, fill=cAll] coordinates {(71.1,6.9)};
\node[lbl, anchor=south east, text=cAll] at (axis cs:71.8,7.8) {\textbf{EBG-all}};
\end{axis}
\end{tikzpicture}

\vspace{1mm}
\caption[Accuracy and harm are not improved together]{%
\textbf{Accuracy and harm do not move together.} Pooled over the four primary models and four decodings,
$n=4288$ runs per condition. \textbf{(a)} Task success and harmful-answer rate for all twelve conditions, both
measured from zero on a common scale. Enforcing record-derived labels (\cond{RT}) roughly triples the harm rate
relative to plain memory (22.6\% against 7.0\%) without improving success. \textbf{(b)} The same conditions placed
by success and harm. \cond{EBG-unsafe} and \cond{EBG-noR3} sit within half a point of \cond{EBG} and are drawn in
the same green but left unlabelled to keep the cluster legible; their exact rates are in
\Cref{tab:rates}. The amber arrow is the central counter-result of this paper: removing the evidence gate, so
that \emph{every} arbitration is enforced rather than only strong-strength ones, moves the system to
\cond{EBG-all} --- $9.7$ percentage points more task success and $3.8$ more harm, with arbitrary picks rising from
13 runs to 166 and false resolution from 1.9\% to 12.9\%. Gating is a trade, not a dominance.\\
\emph{Accessible description:} in panel (b) the evidence-basis conditions occupy the lower right (high success,
low harm); the enforced and LLM-arbiter conditions occupy the upper left (no better success, three to seven times
the harm); an arrow runs up and to the right from \cond{EBG} to \cond{EBG-all}.\\
\emph{Source:} \srcfile{PHASE3/results/per_model_config.csv}; plotted from
\srcfile{figures/data/fig6_accuracy_harm.dat}.}
\label{fig:tradeoff}
\end{figure}


\subsection{Success and harm by condition}

\Cref{fig:tradeoff} places all twelve conditions. Two features of it matter as much as the contrasts above, and
neither favours a simple reading.

\paragraph{The development finding replicates on sealed data.} Metadata-derived labels made agents \emph{worse}
than no labels at all --- \cond{G} 43.5\% against \cond{B} 46.0\%, with more harm (7.9\% against 7.0\%) and more
false conflicts. Enforcing record-derived labels was worse still and far more harmful: \cond{RT} reached 44.2\%
success at 22.6\% harm, roughly tripling the harm rate relative to plain memory while not improving success.
Enforcement is not neutral; applied to labels the evidence does not support, it amplifies them.

\paragraph{Ungated enforcement of the \emph{same} labels scored higher than \cond{EBG}.} \cond{EBG-all} reached
71.1\% success against \cond{EBG}'s 61.3\% --- while more than doubling the harmful rate (6.9\% against 3.1\%),
restoring arbitrary picks from 13 to 166, and raising false resolution from 1.9\% to 12.9\%.
\Cref{sec:ablations} takes this apart. It is the central tension in our results and we do not resolve it in the
mechanism's favour.

% Figure 7 -- per-model effects. Data: figures/data/fig7_ebg_rt.dat (and the three other contrasts).
% Generated from PHASE3/ANALYSIS/CONFIRMATORY_primary_cells.json by figures_source/make_figure_data.py.
\begin{figure}[t]
\centering
\begin{tikzpicture}
\begin{groupplot}[
  paperaxis,
  group style={group size=2 by 1, horizontal sep=14mm},
  width=0.455\linewidth, height=4.4cm,
  ymin=0.5, ymax=4.5, ytick={4,3,2,1},
  yticklabels={Qwen2.5-3B, Qwen2.5-7B, Qwen3-8B, Llama-3.1-8B},
  yticklabel style={font=\scriptsize},
  xmajorgrids,
]

\nextgroupplot[xlabel={task success, EBG $-$ RT (pp)}, xmin=-3, xmax=38, xtick={0,10,20,30},
               title={(a) success}]
\addplot[draw=cInk, line width=0.5pt, dashed] coordinates {(0,0.5) (0,4.5)};
% the other three pre-registered contrasts, for context
\foreach \f in {fig7_ebg_t, fig7_ebg_lg, fig7_ebg_b} {
  \addplot[only marks, mark=o, mark size=1.1pt, draw=cNeutral] table[col sep=tab, x=diff, y expr={\thisrow{y}+0.22}] {figures/data/\f.dat};
}
\addplot[only marks, mark=*, mark size=1.9pt, draw=cEBG, fill=cEBG,
         error bars/.cd, x dir=both, x explicit,
         error bar style={draw=cEBG, line width=0.8pt}]
  table[col sep=tab, x=diff, y=y,
        x error plus expr={\thisrow{hi}-\thisrow{diff}},
        x error minus expr={\thisrow{diff}-\thisrow{lo}}]
  {figures/data/fig7_ebg_rt.dat};

\nextgroupplot[xlabel={harmful answers, EBG $-$ RT (pp)}, xmin=-38, xmax=6, xtick={-30,-20,-10,0},
               title={(b) harm}]
\addplot[draw=cInk, line width=0.5pt, dashed] coordinates {(0,0.5) (0,4.5)};
\addplot[only marks, mark=square*, mark size=1.7pt, draw=cHarm, fill=cHarm,
         error bars/.cd, x dir=both, x explicit,
         error bar style={draw=cHarm, line width=0.8pt}]
  table[col sep=tab, x=harm, y=y,
        x error plus expr={\thisrow{harm_hi}-\thisrow{harm}},
        x error minus expr={\thisrow{harm}-\thisrow{harm_lo}}]
  {figures/data/fig7_ebg_rt.dat};
\node[lbl, text=cEnf, anchor=east, align=right] at (axis cs:-2,1.62)
  {\scriptsize no harm removed:\\[-1pt]\scriptsize $-0.3$\pp{} \ci{-1.0}{0.4}};
\end{groupplot}
\end{tikzpicture}

\vspace{-1mm}
\caption[Per-model effects]{%
\textbf{The effect holds for every model; the mechanism behind it does not.} Filled markers are the per-model
EBG $-$ RT difference with its 95\% family bootstrap interval; open grey markers in (a) are the same model's
EBG $-$ T, EBG $-$ LG and EBG $-$ B differences, shown for context. Per-model estimates are descriptive: the
confirmatory claims rest on the pooled, family-clustered tests of \Cref{fig:forest}. \textbf{(a)} Success improves
for all four models, but by $6.8$--$10.0$\pp{} for Llama-3.1-8B and Qwen2.5-3B against $25.0$--$26.6$\pp{} for
Qwen2.5-7B and Qwen3-8B. \textbf{(b)} Harm falls by 24--28\pp{} for the three Qwen models but by only
$0.3$\pp{} for Llama-3.1-8B, whose interval includes zero. That model abstains on 44.0\% of EBG runs and fails to
format an answer on 11.1\%, so it produces almost no harmful answers under any condition and there is almost
nothing for the mechanism to remove: its success gain comes from converting abstentions into answers, which is a
different causal route from the one the hypotheses describe.\\
\emph{Accessible description:} two dot-and-interval panels over four models; in (a) all intervals lie right of
zero, in (b) three lie far left of zero and the fourth straddles it.\\
\emph{Source:} \srcfile{PHASE3/ANALYSIS/CONFIRMATORY_primary_cells.json}; plotted from
\srcfile{figures/data/fig7_ebg_rt.dat}.}
\label{fig:permodel}
\end{figure}


\subsection{Per-model effects}

Per-model differences are descriptive; the confirmatory claims rest on the pooled tests. The direction is positive
for every model on every contrast and every per-model interval excludes zero, but the magnitude is not uniform
(\Cref{fig:permodel}): on H-primary the two mid-sized Qwen models gain $26.6$ and $25.0$\pp{} against $10.0$\pp{}
for Qwen2.5-3B and $6.8$\pp{} for Llama-3.1-8B.

The Llama-3.1-8B result deserves emphasis because it was pre-stated as a risk. On the development split this model
showed essentially no benefit ($+1.1$\pp), and the pre-registration said a repeat on eval3 would be a failure to
generalise across families and would be reported as such. On eval3 the benefit is small but positive with an
interval excluding zero. The \emph{harm} reduction, however, does not transfer at all ($-0.3$\pp{}
\ci{-1.0}{0.4}): this model abstains rather than committing --- 44.0\% unnecessary abstention under \cond{EBG},
60.7\% under \cond{B} --- and formats no parsable answer on 11.1\% of runs, so it was producing almost no harmful
answers for the mechanism to remove. Its gain comes from converting abstentions into answers, a different causal
route from the one the hypotheses describe.

\subsection{Secondary cells}

The two greedy-only secondary models are never pooled with the primary ones. Qwen2.5-14B shows the largest
\cond{EBG} $-$ \cond{RT} effect in the study, $+29.9$\pp{} \ci{23.0}{36.9}. Llama-3.2-3B, which failed the
readiness criterion before the campaign, shows $+12.3$\pp{} \ci{6.3}{18.3} from a very low base --- 41.8\% success
under \cond{EBG} against 7.1\% under \cond{B}, with 88.7\% unnecessary abstention under \cond{B} --- so its
differences measure recovery from near-total failure rather than the effect the hypotheses concern.

\subsection{Absolute performance}

No condition is good in absolute terms. The best, \cond{EBG}, answers 61.3\% of pooled primary runs correctly, and
46.0\% of the runs whose tasks have a definite answer rather than a genuine conflict. Qwen2.5-7B and Qwen3-8B reach
75.0\% and 77.8\% under \cond{EBG}; Llama-3.1-8B reaches 33.0\%. The contribution is a relative one about
mechanism design, not a claim that small agents handle conflicting persistent memory well.
